% Part: second-order-logic
% Chapter: sol-and-sets
% Section: comparing-sets

\documentclass[../../../include/open-logic-section]{subfiles}

\begin{document}

\olfileid[tr]{sol}{set}{cmp}

\olsection{Kümelerin Karşılaştırılması}

\begin{prop}
!!{formula} $\lforall[x][(X(x) \lif Y(x))]$ alt küme bağıntısını tanımlar;
yani $\Sat{M}{\lforall[x][(X(x) \lif Y(x))]}[s]$ ancak ve ancak
$s(X)\subseteq s(Y)$ ise.
\end{prop}

\begin{prop}
!!{formula} $\lforall[x][(X(x) \liff Y(x))]$ kümeler üzerindeki özdeşlik
bağıntısını tanımlar; yani
$\Sat{M}{\lforall[x][(X(x) \liff Y(x))]}[s]$ ancak ve ancak $s(X)=s(Y)$ ise.
\end{prop}

\begin{prop}
!!{formula} $\lexists[x][X(x)]$ boş olmama özelliğini tanımlar; yani
$\Sat{M}{\lexists[x][X(x)]}[s]$ ancak ve ancak $s(X)\neq\emptyset$ ise.
\end{prop}

Bir $X$ kümesi, bir $Y$ kümesinden daha büyük değildir,
$\cardle{X}{Y}$, ancak ve ancak !!a{injective} $f\colon X\to Y$ fonksiyonu
varsa. Bir fonksiyonun bire bir olduğunu ve $X$ içindeki argümanlardaki
değerlerinin $Y$ içinde bulunduğunu ifade edebildiğimizden, tanım kümesinin
alt kümeleri üzerinde daha büyük olmama bağıntısını da tanımlayabiliriz.

\begin{prop}
\[
\lexists[u][(\lforall[x][(X(x) \lif Y(u(x)))] \land \lforall[x][\lforall[y][(\eq[u(x)][u(y)] \lif
      \eq[x][y])]])]
\]
formülü daha büyük olmama bağıntısını tanımlar.
\end{prop}

İki küme aynı büyüklüktedir ya da ``eş güçlüdür'', $\cardeq{X}{Y}$, ancak ve
ancak !!a{bijective} $f\colon X\to Y$ fonksiyonu varsa.

\begin{prop}
\begin{multline*}
  \lexists[u][(\lforall[x][(X(x) \lif Y(u(x)))] \land {}\\
    \lforall[x][\lforall[y][(\eq[u(x)][u(y)] \lif \eq[x][y])]]
  \land {} \\\\lforall[y][(Y(y) \lif \lexists[x][(X(x)
      \land \eq[y][u(x)])])])]
\end{multline*}
formülü eş güçlü olma bağıntısını tanımlar.
\end{prop}

Bu !!{formula}s{}i sırasıyla $X\subseteq Y$, $X=Y$, $X\neq\emptyset$,
$\cardle{X}{Y}$ ve $\cardeq{X}{Y}$ ile kısaltacağız. (Kümeler $X$ ve $Y$
hakkında biçimsel olmayan biçimde konuşurken de aynı gösterimi kullandığımızdan
bu biraz kafa karıştırabilir; fakat burada gösterim, $X$ ve $Y$ tek konumlu
bağıntı değişkenlerini içeren ikinci dereceden mantık !!{formula}s{}inin
kısaltmasıdır.)

\begin{prop}
!!{sentence} $\lforall[X][\lforall[Y][((\cardle{X}{Y} \land
    \cardle{Y}{X}) \lif \cardeq{X}{Y})]]$ geçerlidir.
\end{prop}

\begin{proof}
Bu !!{sentence}, her $X\subseteq\Domain{M}$ ve $Y\subseteq\Domain{M}$ alt
kümeleri için $\cardle{X}{Y}$ ve $\cardle{Y}{X}$ olduğunda
$\cardeq{X}{Y}$ oluyorsa !!a{structure} $\Struct{M}$ içinde sağlanır. Bu
ise \emph{her} $X$ ve $Y$ kümesi için geçerlidir: Schröder--Bernstein
Teoremidir.
\end{proof}


\end{document}
